\section*{[2023 PREP EXAM] Problem 5: Ritz projection, error estimates for the heat equation}
We consider the semi-discrete (in space) heat equation: find
\(u_h \in C^1([0,T];V_h)\) such that
\[
    \langle u_{h,t}(t),v_h\rangle + \langle \nabla u_h(t),\nabla v_h\rangle
    = \langle f(t),v_h\rangle \quad \forall v_h\in V_h,\ \forall t\in(0,T), \tag{\text{semi}}
\]
with initial condition
\[
    u_h(0,x) = \Pi_{V_h} u_0(x), \tag{\text{semi-init}}
\]
where \(\Pi_{V_h}:H^1(\Omega)\to V_h\) is the \(L^2\)-projection and
\(V_h\subset H^1_0(\Omega)\).
We denote \((\cdot,\cdot) := \langle\cdot,\cdot\rangle\) the \(L^2(\Omega)\) inner product
and use the shorthand
\[
    \|v\| := \|v\|_{L^2(\Omega)},\qquad |v|_1 := \|\nabla v\|_{L^2(\Omega)}.
\]
The elliptic bilinear form is
\[
    a(u,v) := (\nabla u,\nabla v) \quad \text{for } u,v\in H^1_0(\Omega).
\]
The \emph{Ritz projection} \(R_h:H^1_0(\Omega)\to V_h\) is defined by
\[
    a(R_h v - v,\chi) = 0 \quad \forall \chi\in V_h. \tag{\text{Ritz-def}}
\]
We are also given a regularity/approximation assumption:
for a family of regular triangulations \(\{T_h\}\) and corresponding FE spaces
\(V_h\subset H^1_0(\Omega)\), there exist \(r\ge 2\) and \(C>0\) such that
for \(1\le s\le r\) and small enough \(h\),
\[
    \inf_{\chi\in V_h} \bigl(\|u_0 - \chi\| + h\,|\nabla(u_0 - \chi)|_0\bigr)
    \;\le\; C h^s \|u_0\|_{H^s(\Omega)},
    \qquad u_0\in H^s(\Omega)\cap H^1_0(\Omega). \tag{\text{reg-ass}}
\]
\subsection*{(a) Stability of the Ritz projection in the $H^1$ semi-norm}
We want to prove
\[
    |R_h v|_1 \le |v|_1 \quad \forall v\in H^1_0(\Omega).
\]
Let \(v\in H^1_0(\Omega)\) and set \(w_h := R_h v \in V_h\).
By the definition (reg-ass), choosing \(\chi = w_h\) gives
\[
    a(R_h v - v, w_h) = 0
    \quad\Longrightarrow\quad
    a(w_h,w_h) = a(v,w_h).
\]
Expanding \(a\),
\[
    (\nabla w_h,\nabla w_h) = (\nabla v,\nabla w_h).
\]
The left-hand side is \(|w_h|_1^2\), so
\[
    |w_h|_1^2 = (\nabla v,\nabla w_h)
    \le |\nabla v|_0 |\nabla w_h|_0
    = |v|_1\,|w_h|_1,
\]
where we used the Cauchy--Schwarz inequality.
If \(|w_h|_1 = 0\) there is nothing to prove (then \(R_h v\equiv 0\), so
\(|R_h v|_1 = 0 \le |v|_1\)).
Otherwise, divide by \(|w_h|_1>0\):
\[
    |w_h|_1 \le |v|_1.
\]
Thus
\[
    \boxed{|R_h v|_1 \le |v|_1 \quad \forall v\in H^1_0(\Omega),}
\]
i.e.\ the Ritz projection is stable in the \(H^1\) semi-norm.
\subsection*{(b) Optimal $H^1$-semi-norm convergence of $R_h$}
We want to show that for smooth enough \(v\in H^s(\Omega)\cap H^1_0(\Omega)\),
\(1\le s\le r\), the Ritz projection converges \emph{optimally} in the
\(H^1\)-semi-norm:
\[
    |v - R_h v|_1 \le C h^{s-1} \|v\|_{H^s(\Omega)}.
\]
\textbf{Best-approximation property.}
Let \(v\in H^1_0(\Omega)\), and let \(e := v - R_h v\).
By (ritz-def), for all \(\chi\in V_h\),
\[
    a(e,\chi) = a(v - R_h v, \chi) = 0.
\]
In particular,
\[
    a(e,e) = a(e, v - R_h v)
    = a(e, v - \chi) + a(e,\chi - R_h v).
\]
But \(\chi - R_h v \in V_h\), so \(a(e,\chi - R_h v)=0\). Hence
\[
    a(e,e) = a(e, v - \chi) \quad \forall \chi\in V_h.
\]
Using Cauchy--Schwarz for \(a(\cdot,\cdot)\) in the \(H^1\)-semi-norm:
\[
    |e|_1^2 = a(e,e) = a(e, v - \chi)
    \le |e|_1\,|v - \chi|_1.
\]
If \(e\neq 0\), divide by \(|e|_1\):
\[
    |v - R_h v|_1 = |e|_1 \le |v - \chi|_1 \quad \forall \chi\in V_h.
\]
Therefore,
\[
    |v - R_h v|_1 \le \inf_{\chi\in V_h} |v - \chi|_1. \tag{\text{best-approx}}
\]
\textbf{Use of the approximation/regularity assumption.}
Let \(v\in H^s(\Omega)\cap H^1_0(\Omega)\), with \(1\le s\le r\).
By the assumption (reg-ass), for each such \(v\) there exists some
\(\chi_h\in V_h\) (approximately attaining the infimum) such that
\[
    \|v - \chi_h\| + h\,|\nabla(v - \chi_h)|_0
    \;\le\; C h^s \|v\|_{H^s(\Omega)}.
\]
In particular,
\[
    h\,|v - \chi_h|_1
    = h \|\nabla(v - \chi_h)\|_0
    \le C h^s \|v\|_{H^s(\Omega)},
\]
so
\[
    |v - \chi_h|_1 \le C h^{s-1} \|v\|_{H^s(\Omega)}.
\]
Since
\[
    \inf_{\chi\in V_h} |v - \chi|_1 \le |v - \chi_h|_1,
\]
combining with (best-approx) yields
\[
    |v - R_h v|_1
    \le \inf_{\chi\in V_h} |v - \chi|_1
    \le C h^{s-1} \|v\|_{H^s(\Omega)}.
\]
Thus, for \(1\le s\le r\),
\[
    \boxed{
        |v - R_h v|_1 \le C h^{s-1} \|v\|_{H^s(\Omega)}
    }
\]
which is the \emph{optimal} order in \(H^1\)-semi-norm (matching the standard FE interpolation
order).
\subsection*{(c) Parabolic error splitting and an $L^2$-error estimate}
We split the total error into elliptic and parabolic parts:
\[
    u_h(t) - u(t)
    = \underbrace{u_h(t) - R_h u(t)}_{\theta(t)}
    + \underbrace{R_h u(t) - u(t)}_{\rho(t)}
    =: \theta(t) + \rho(t).
\]
We will prove the first of the two proposed estimates:
\emph{Claim:} Let \(u\in C^1([0,T];H^s(\Omega))\), \(1\le s\le r\), solve the heat equation
\[
    (u_t,v) + a(u,v) = (f,v)\quad \forall v\in H^1_0(\Omega),
\]
and let \(u_h\) be the semi-discrete solution (semi)--(semi-init).
Assume the regularity/approximation condition (reg-ass) and that
the Ritz projection satisfies the optimal \(L^2\) estimate
\[
    \|w - R_h w\| \le C h^s \|w\|_{H^s(\Omega)}
    \quad\text{for } w\in H^s(\Omega)\cap H^1_0(\Omega),\ 1\le s\le r. \tag{\text{Ritz-L2}}
\]
Then, for all \(t\in[0,T]\),
\[
    \|u_h(t) - u(t)\|
    \le \|\Pi_{V_h}u_0 - u_0\|
    + C h^s\Bigl(\|u_0\|_{H^s(\Omega)}
    + \int_0^t \|u_\tau(\tau)\|_{H^s(\Omega)}\,d\tau\Bigr). \tag{\text{L2-final}}
\]
\textbf{Step 1: Error equation for $\theta$.}
The continuous weak form is:
\[
    (u_t(t),v) + a(u(t),v) = (f(t),v) \quad \forall v\in H^1_0(\Omega).
\]
The semi-discrete form is:
\[
    (u_{h,t}(t),v_h) + a(u_h(t),v_h) = (f(t),v_h) \quad \forall v_h\in V_h.
\]
Subtracting and restricting to test functions \(v_h\in V_h\), we obtain
\[
    (u_{h,t} - u_t, v_h) + a(u_h - u, v_h) = 0
    \quad \forall v_h\in V_h.
\]
Write \(u_h - u = \theta + \rho\) and \(u_{h,t} - u_t = \theta_t + \rho_t\):
\[
    (\theta_t + \rho_t, v_h) + a(\theta + \rho, v_h) = 0
    \quad\forall v_h\in V_h.
\]
By definition of the Ritz projection,
\[
    a(\rho(t), v_h) = a(R_h u(t) - u(t),v_h) = 0
    \quad \forall v_h\in V_h,\ \forall t,
\]
so $a(\rho,v_h)$ vanishes. Hence
\[
    (\theta_t,v_h) + (\rho_t,v_h) + a(\theta,v_h) = 0
    \quad\forall v_h\in V_h.
\]
Take \(v_h = \theta(t)\in V_h\):
\[
    (\theta_t,\theta) + (\rho_t,\theta) + a(\theta,\theta) = 0.
\]
Note
\[
    (\theta_t,\theta) = \frac{1}{2}\frac{d}{dt}\|\theta(t)\|^2,
    \qquad
    a(\theta,\theta) = |\theta(t)|_1^2.
\]
Thus
\[
    \frac{1}{2}\frac{d}{dt}\|\theta(t)\|^2 + |\theta(t)|_1^2
    = -(\rho_t(t),\theta(t)). \tag{\text{theta-energy}}
\]
\textbf{Step 2: Differential inequality for $\|\theta(t)\|$.}
Using Cauchy--Schwarz in (theta-energy):
\[
    |(\rho_t(t),\theta(t))| \le \|\rho_t(t)\|\,\|\theta(t)\|.
\]
Dropping the nonnegative term $|\theta(t)|_1^2$ gives
\[
    \frac{1}{2}\frac{d}{dt}\|\theta(t)\|^2
    \le \|\rho_t(t)\|\,\|\theta(t)\|.
\]
When $\|\theta(t)\|\neq 0$, we can divide by $\|\theta(t)\|$ and use
\[
    \frac{d}{dt}\|\theta(t)\|
    = \frac{1}{\|\theta(t)\|}\frac{d}{dt}\Bigl(\frac12\|\theta(t)\|^2\Bigr),
\]
to obtain
\[
    \frac{d}{dt}\|\theta(t)\| \le \|\rho_t(t)\|.
\]
Integrating from $0$ to $t$:
\[
    \|\theta(t)\| \le \|\theta(0)\| + \int_0^t \|\rho_\tau(\tau)\|\,d\tau. \tag{\text{theta-bound}}
\]
\textbf{Step 3: Initial value $\theta(0)$ and bounds for $\rho$.}
By definition,
\[
    \theta(0) = u_h(0) - R_h u(0) = \Pi_{V_h}u_0 - R_h u_0.
\]
Hence
\[
    \|\theta(0)\|
    \le \|\Pi_{V_h}u_0 - u_0\| + \|u_0 - R_h u_0\|
    = \|\Pi_{V_h}u_0 - u_0\| + \|\rho(0)\|.
\]
Next, recall \(\rho(t) = R_h u(t) - u(t)\) and
\(\rho_t(t) = R_h u_t(t) - u_t(t)\).
Using the \(L^2\)-estimate (Ritz-L2), for each \(t\in[0,T]\),
\[
    \|\rho(t)\| = \|R_h u(t) - u(t)\|
    \le C h^s \|u(t)\|_{H^s(\Omega)},
\]
\[
    \|\rho_t(t)\| = \|R_h u_t(t) - u_t(t)\|
    \le C h^s \|u_t(t)\|_{H^s(\Omega)}.
\]
In particular,
\[
    \|\rho(0)\| \le C h^s \|u_0\|_{H^s(\Omega)}.
\]
\textbf{Step 4: Bound $\|u(t)\|_{H^s}$ in terms of $\|u_0\|_{H^s}$ and $u_t$.}
For each $t\in[0,T]$, by the fundamental theorem of calculus in $H^s(\Omega)$,
\[
    u(t) = u_0 + \int_0^t u_\tau(\tau)\,d\tau,
\]
hence
\[
    \|u(t)\|_{H^s}
    \le \|u_0\|_{H^s} + \int_0^t \|u_\tau(\tau)\|_{H^s}\,d\tau.
\]
Therefore,
\[
    \|\rho(t)\|
    \le C h^s\Bigl(\|u_0\|_{H^s} + \int_0^t \|u_\tau(\tau)\|_{H^s}\,d\tau\Bigr),
\]
and
\[
    \int_0^t \|\rho_\tau(\tau)\|\,d\tau
    \le C h^s \int_0^t \|u_\tau(\tau)\|_{H^s}\,d\tau.
\]
\textbf{Step 5: Combine bounds for the total error.}
Recall
\[
    e(t) := u_h(t) - u(t) = \theta(t) + \rho(t),
\]
so
\[
    \|e(t)\|\le \|\theta(t)\| + \|\rho(t)\|.
\]
Using (theta-bound) and the bounds above,
\[
    \begin{aligned}
        \|\theta(t)\|
         & \le \|\Pi_{V_h}u_0 - u_0\| + \|\rho(0)\|
        + \int_0^t \|\rho_\tau(\tau)\|\,d\tau       \\
         & \le \|\Pi_{V_h}u_0 - u_0\|
        + C h^s \|u_0\|_{H^s}
        + C h^s \int_0^t \|u_\tau(\tau)\|_{H^s}\,d\tau.
    \end{aligned}
\]
Thus
\[
    \begin{aligned}
        \|u_h(t) - u(t)\|
         & \le \|\theta(t)\| + \|\rho(t)\|             \\
         & \le \|\Pi_{V_h}u_0 - u_0\|
        + C h^s \|u_0\|_{H^s}
        + C h^s \int_0^t \|u_\tau(\tau)\|_{H^s}\,d\tau \\
         & \quad
        + C h^s \Bigl(\|u_0\|_{H^s} + \int_0^t \|u_\tau(\tau)\|_{H^s}\,d\tau\Bigr).
    \end{aligned}
\]
Absorbing constants (and the two identical \(h^s\)-terms) into a new \(C\), we obtain
\[
    \|u_h(t) - u(t)\|
    \le \|\Pi_{V_h}u_0 - u_0\|
    + C h^s \Bigl(\|u_0\|_{H^s} + \int_0^t \|u_\tau(\tau)\|_{H^s}\,d\tau\Bigr),
\]
which is exactly (L2-final).
Thus we have:
\[
    \boxed{
        \|u_h(t) - u(t)\|
        \le \|\Pi_{V_h}u_0 - u_0\|
        + C h^s\Bigl(\|u_0\|_{H^s(\Omega)} + \int_0^t \|u_\tau(\tau)\|_{H^s(\Omega)}\,d\tau\Bigr)
    }
\]
for all \(t\in[0,T]\), proving the first of the requested parabolic a priori error estimates.

